\documentclass[11pt,a4paper]{article}
\usepackage{amsmath,amssymb,amsthm}
\usepackage[margin=2.5cm]{geometry}
\usepackage{hyperref}

\newtheorem{definition}{Definition}[section]
\newtheorem{theorem}[definition]{Theorem}
\newtheorem{proposition}[definition]{Proposition}
\newtheorem{remark}[definition]{Remark}
\newtheorem{conjecture}[definition]{Conjecture}

\title{Hyperseed Formalization:\\
  NFTs and the Gamification of the Economy}
\author{ProtomegaTron (automated formalization)}
\date{2026-09-18}

\begin{document}
\maketitle

\begin{abstract}
We formalize Ben Goertzel's ``NFTs and the Gamification of the Economy'' (2022-03-07)
within the Hyperseed ontology.
\end{abstract}

\section{Fiber play: gamified economics}

\begin{proposition}[Gamification = fiber play --- claims 1, 5, 6]
Gamified economics is fiber play --- exploring fiber structure beyond
base-space survival:
\[
  \text{When } B_{\text{material}} \text{ is abundant} \implies
  \text{activity} \to F_{\text{play}}
\]
Economic activity becomes exploration of fiber structure (social position,
creative expression, status games) rather than base-space resource acquisition.
\end{proposition}

\section{NFTs as fiber tokens}

\begin{proposition}[NFTs = fiber tokens --- claims 2--4, 7]
NFTs are markers of position in social fiber space:
\[
  \text{NFT}(a) \in F_{\text{social}}(b_a) \quad \text{marking ownership, membership, status}
\]
The token itself is a fiber element --- its value comes from its position
in the social fiber bundle, not from its base-space (material) properties.
\end{proposition}

\section{Post-scarcity as thick base}

\begin{proposition}[Post-scarcity = thick base --- claims 5, 8]
When base-space resources become abundant (thick base), activity shifts
to fiber exploration:
\[
  \lim_{|B_{\text{resources}}| \to \infty} \text{focus} \to F_{\text{meaning}}
\]
Post-scarcity economics is economics where the fiber is the point, not
the base.
\end{proposition}
\end{document}
